% 릴레이식 유닛 모듈의 개략 치수. 단위 mm. 그림 fig:hw-relay-dim-unit
% 근거: hardware/relay/layout.md 3.1-3.5절. 카드 250 x 228, 커넥터 자리 h1 = 12-82, h2 = 86-156 (gen_shifter_pcb.py).
\begin{tikzpicture}[x=0.02cm,y=0.02cm,
    dim/.style={<->,cgray,line width=0.4pt},
    dl/.style={font=\tiny,text=black,inner sep=1pt},
    lab/.style={font=\tiny,align=center}]
  % ================= (가) 앞 =================
  \node[font=\footnotesize\bfseries,anchor=west] at (-10,372) {(가) 앞에서 본 유닛 한 자리};
  % 자리(열 피치 240, 단 피치 290)
  \draw[draw=cgray,dashed] (0,0) rectangle (240,290);
  % 옆판
  \draw[draw=cgray,fill=clight] (3,0) rectangle (5,290);
  \draw[draw=cgray,fill=clight] (235,0) rectangle (237,290);
  % 카드 (a4 시프터 보드)
  \draw[draw=cblue,fill=cbluebg,line width=0.6pt] (6,20) rectangle (234,270);
  % 릴레이 26 x 5단
  \foreach \b in {0,...,4} {
    \foreach \i in {0,...,25} {
      \draw[draw=cgray,fill=white,line width=0.2pt] ({16.5+\i*8},{58+\b*38-7}) rectangle ++(7,15);
    }
    \fill[cgreen] (229,{58+\b*38+7}) circle (2.2);
  }
  \node[dl,anchor=east] at (-4,58) {1칸 단};
  \node[dl,anchor=east] at (-4,210) {16칸 단};
  \draw[cgray,line width=0.3pt] (-3,58) -- (14,58);
  \draw[cgray,line width=0.3pt] (-3,210) -- (14,210);
  % 아래 변: J1(h1), J2(h2), J5
  \foreach \a/\b/\t in {18/88/J1, 92/162/J2} {
    \draw[draw=cgray,fill=clight] (\a,20) rectangle (\b,32);
    \draw[draw=corange,fill=corangebg] (\a,5) rectangle (\b,20);
    \node[dl] at ({(\a+\b)/2},12) {\t};
  }
  \draw[draw=cgray,fill=clight] (166,20) rectangle (198,32);
  \draw[draw=cpurple,fill=white] (166,5) rectangle (198,20);
  \node[dl] at (182,12) {J5};
  % 위 변: J3(h1), J4(h2)
  \foreach \a/\b/\t in {18/88/J3, 92/162/J4} {
    \draw[draw=cgray,fill=clight] (\a,258) rectangle (\b,270);
    \draw[draw=corange,fill=corangebg] (\a,270) rectangle (\b,285);
    \node[dl] at ({(\a+\b)/2},277) {\t};
  }
  \node[dl,text=cgray] at (53,300) {h1};
  \node[dl,text=cgray] at (127,300) {h2};
  % 구멍
  \foreach \x/\y in {11/25, 229/25, 11/265, 229/265} { \draw[cgray] (\x,\y) circle (1.8); }
  % 치수
  \draw[cgray,line width=0.3pt] (6,272) -- (6,322);
  \draw[cgray,line width=0.3pt] (234,272) -- (234,322);
  \draw[dim] (6,314) -- (234,314) node[midway,dl,fill=white] {카드 228};
  \draw[cgray,line width=0.3pt] (0,-2) -- (0,-26);
  \draw[cgray,line width=0.3pt] (240,-2) -- (240,-26);
  \draw[dim] (0,-18) -- (240,-18) node[midway,dl,fill=white] {열 피치 240};
  \foreach \y in {0,20,270,290} { \draw[cgray,line width=0.3pt] (242,\y) -- (300,\y); }
  \draw[dim] (262,20) -- (262,270) node[midway,dl,fill=white,rotate=90] {카드 250};
  \draw[dim] (262,0) -- (262,20);
  \draw[dim] (262,270) -- (262,290);
  \node[dl,anchor=west] at (265,10) {20};
  \node[dl,anchor=west] at (265,280) {20};
  \draw[dim] (292,0) -- (292,290) node[midway,dl,fill=white,rotate=90] {단 피치 290};

  % ================= (나) 옆 =================
  \begin{scope}[shift={(400,0)}]
    \node[font=\footnotesize\bfseries,anchor=west] at (-40,372) {(나) 옆에서 본 겹침 (왼쪽이 관람객)};
    % 덮개, 장착판, 뒤 덮개
    \draw[draw=cblue,fill=cbluebg] (-30,0) rectangle (-25,290);
    \draw[draw=cgray,fill=clight] (227,0) rectangle (230,290);
    \draw[draw=cgray,fill=clight] (290,0) rectangle (293,290);
    \node[lab,text=cgray,rotate=90] at (260,145) {배선 공간};
    % 카드: 부품면의 방향을 옅은 띠로
    \draw[draw=none,fill=cbluebg] (-10,20) rectangle (0,270);
    \foreach \z in {25,50,75,100,125,150,175} { \draw[draw=none,fill=clight] (\z+2,20) rectangle (\z+12,270); }
    \draw[draw=none,fill=cyellowbg] (205,20) rectangle (215,270);
    \foreach \z in {0,75,125,175} { \draw[draw=cblue,fill=cblue] (\z,20) rectangle (\z+2,270); }
    \foreach \z in {50,100,150} { \draw[draw=corange,fill=corange] (\z,20) rectangle (\z+2,270); }
    \draw[draw=cred,fill=cred] (25,20) rectangle (27,270);
    \draw[draw=cgreen,fill=cgreen] (215,20) rectangle (217,270);
    \foreach \z/\t in {0/{S4 ($a_4$)}, 25/U, 50/{$W_L$}, 75/{S3 ($a_3$)}, 100/{$W_M$}, 125/{S2 ($a_2$)},
                       150/{$W_R$}, 175/{S1 ($a_1$)}, 215/C} {
      \node[dl,rotate=90,anchor=west] at (\z+1,294) {\t};
    }
    % 점퍼: 아래 h1, 아래 h2, 위 h1, 위 h2
    \foreach \a/\b in {0/25, 50/75, 100/125, 150/175} {
      \draw[corange,line width=1pt] (\a+1,20) -- (\a+1,13) -- (\b+1,13) -- (\b+1,20);
      \draw[corange,line width=1pt] (\a+1,270) -- (\a+1,283) -- (\b+1,283) -- (\b+1,270);
    }
    \foreach \a/\b in {0/50, 75/100, 125/150} {
      \draw[corange,line width=1pt] (\a+1,20) -- (\a+1,6) -- (\b+1,6) -- (\b+1,20);
      \draw[corange,line width=1pt] (\a+1,270) -- (\a+1,277) -- (\b+1,277) -- (\b+1,270);
    }
    % 리드: IN (위 h1), OUT (아래 h2)
    \draw[->,cblue,line width=0.8pt] (255,277) -- (176,277) -- (176,270);
    \node[dl,text=cblue,anchor=south] at (240,279) {\code{IN}};
    \draw[->,cblue,line width=0.8pt] (176,20) -- (176,6) -- (255,6);
    \node[dl,text=cblue,anchor=north] at (240,4) {\code{OUT}};
    % 치수
    \foreach \z in {-30,0,25,175,215,230,293,300} { \draw[cgray,line width=0.3pt] (\z,-2) -- (\z,-26); }
    \draw[dim] (-30,-18) -- (0,-18) node[midway,dl,fill=white] {30};
    \draw[dim] (0,-18) -- (25,-18) node[midway,dl,fill=white] {25};
    \draw[dim] (25,-18) -- (175,-18) node[midway,dl,fill=white] {6 × 25};
    \draw[dim] (175,-18) -- (215,-18) node[midway,dl,fill=white] {40};
    \draw[dim] (230,-18) -- (293,-18) node[midway,dl,fill=white] {63};
    \draw[cgray,line width=0.3pt] (-30,-26) -- (-30,-56);
    \draw[cgray,line width=0.3pt] (300,-26) -- (300,-56);
    \draw[dim] (-30,-48) -- (300,-48) node[midway,dl,fill=white] {뱅크의 깊이 330};
    % 범례
    \node[lab,anchor=west,text=cblue] at (-30,-78) {파랑: 시프터 보드};
    \node[lab,anchor=west,text=corange] at (95,-78) {주황: 로터 카드와 점퍼};
    \node[lab,anchor=west,text=cred] at (-30,-100) {빨강: 반사판 카드};
    \node[lab,anchor=west,text=cgreen] at (95,-100) {초록: 계수기 유닛 보드};
    \node[lab,anchor=west,text=cgray] at (230,-78) {옅은 띠:};
    \node[lab,anchor=west,text=cgray] at (230,-100) {부품이 있는 쪽};
  \end{scope}
\end{tikzpicture}
